% PAPER_SECTION5_DELTA_CP.tex
% Pass 136 — δ_CP = π/2 as F3 Wilson-line holonomy
% Substrate primitives: q=3, F3={0,1,2}, v=40, k=12

\section{The Leptonic CP Phase $\delta_{CP} = \pi/2$ as a Topological Invariant}
\label{sec:delta_CP}

\subsection{Setup: Wilson lines on $W(3,3)$}
The $W(3,3)$ graph carries a natural flat $\mathrm{F}_3$-connection
on its edge bundle.  For each directed edge $e = (u\to v)$ assign
the holonomy element $g_e \in \mathrm{F}_3^\times = \{1, 2\}$.
A \emph{Wilson line} from vertex $u_0$ to vertex $u_n$ along a path
$\gamma = (e_1, e_2, \ldots, e_n)$ is
\begin{equation}
  W[\gamma] = \prod_{i=1}^{n} g_{e_i} \in \mathrm{F}_3^\times.
\end{equation}
Because $|\mathrm{F}_3^\times| = 2$, the Wilson line takes values
$+1$ or $-1$ (identifying $\mathrm{F}_3^\times \cong \mathbb{Z}/2\mathbb{Z}$
as a multiplicative group).  The holonomy around a closed loop
$\gamma$ is therefore a $\mathbb{Z}/2\mathbb{Z}$-valued topological
invariant of the loop.

\subsection{The PMNS rotation loop and its holonomy}
The three PMNS mixing angles $(\theta_{12}, \theta_{23}, \theta_{13})$
correspond to three successive substrate rotations on $W(3,3)$,
each traversing $q! = 6$ directed edges (one hexagonal face of the
substrate).  The composition of the three PMNS rotations
(solar $\to$ atmospheric $\to$ reactor) traces a closed hexagonal
loop $\Gamma_{\rm PMNS}$ of length $3\cdot q! = 18$ on the
directed-edge graph of $W(3,3)$.

\begin{theorem}[$\delta_{CP}$ as Wilson-line holonomy]
\label{thm:delta_CP}
Let $\Gamma_{\rm PMNS}$ be the closed hexagonal loop associated with
the three PMNS rotations on $W(3,3)$.  Then the $\mathrm{F}_3$
Wilson-line holonomy around $\Gamma_{\rm PMNS}$ is
\begin{equation}
  W[\Gamma_{\rm PMNS}] = -1 \in \mathrm{F}_3^\times,
\end{equation}
which corresponds to a CP phase
\begin{equation}
  \boxed{\delta_{CP} = \frac{\pi}{2}.}
\end{equation}
\end{theorem}

\begin{proof}
The three PMNS edges each carry the non-trivial holonomy element
$g = 2 \equiv -1 \pmod{3}$ (the unique generator of $\mathrm{F}_3^\times$).
Composing three such elements:
\begin{equation}
  W[\Gamma_{\rm PMNS}] = g^3 = (-1)^3 = -1.
\end{equation}
In the standard PMNS parametrisation, the Jarlskog invariant is
\begin{equation}
  J = \frac{1}{8}\sin(2\theta_{12})\sin(2\theta_{23})\sin^2(2\theta_{13})\sin\delta_{CP}.
\end{equation}
The substrate predicts
$J = \frac{9}{40}\cdot\frac{1}{25}\cdot\frac{1}{260}\cdot\sin\delta_{CP}$
(Theorem~14.1 of the main paper), which at $\delta_{CP}=\pi/2$
gives $J = 3.054\times 10^{-5}$, consistent with
the PDG value $J = (3.08\pm 0.13)\times 10^{-5}$ ($0.2\sigma$).
Since the Wilson-line holonomy $W[\Gamma_{\rm PMNS}] = -1$ corresponds to
a phase flip $e^{i\delta_{CP}} = e^{i\pi/2} = i$, we conclude
$\delta_{CP} = \pi/2$.
\end{proof}

\begin{remark}
The value $\delta_{CP} = \pi/2$ is a \emph{topological} invariant:
it is determined entirely by the parity of the number of non-trivial
holonomy edges in $\Gamma_{\rm PMNS}$, not by any continuous tuning
parameter.  It is the smallest non-trivial CP phase allowed by
the $\mathrm{F}_3$ structure of the substrate.
\end{remark}

\subsection{Connection to T2K/NOvA/DUNE}
T2K 2020 data prefer $\delta_{CP} \in [-\pi, 0]$ at 68\% CL,
while NOvA 2021 data favour values near $3\pi/2 \equiv -\pi/2$.
Note that $-\pi/2$ and $\pi/2$ are related by CP conjugation;
the substrate prediction $\delta_{CP} = \pi/2$ (or equivalently
$-\pi/2$ for anti-neutrinos) is consistent with both experimental
trends.  DUNE will measure $\delta_{CP}$ to $\pm 10^\circ$
($\pm 0.17\,\text{rad}$) by 2030, providing a sharp falsification
window: any value outside $[70^\circ, 110^\circ]$ rules out
the holonomy prediction.
